\section*{Butter Chicken aus dem Ofen}
\ihead{}\chead{Personen:2}\ohead{}
\ifoot{Vorbereitungszeit:30 min}\cfoot{}\ofoot{Kochzeit:30 min}
{\Large Zutaten}
\begin{multicols}{2}
\begin{itemize}
\item \SI{500}{g} Hähnchenbrust
\item \SI{1}{TL} Paprikapulver
\item \SI{1}{EL} Limonensaft oder Zitronensaft
\item \SI{1}{TL} Salz
\item \SI{150}{g} Joghurt
\item \SI{1}{TL} Cayennepfeffer
\item \SI{1}{EL} Garam Masala
\item \num{1} Stück Ingwer, daumengroß
\item \num{1} Knoblauchzehe
\item \SI{4}{EL} Butter
\item \num{1} Zwiebel
\item \SI{500}{g} passierte Tomaten
\item \SI{1}{TL} Zimt
\item \SI{1}{TL} Salz
\item \SI{2}{TL} Cayennepfeffer
\item \num{1} Stück Ingwer, daumengroß
\item \num{1} Knoblauchzehe
\item \SI{1}{EL} Honig
\item \SI{150}{ml} Sahne
\item einige Korianderblätter, optional
\end{itemize}
% \columnbreak
\end{multicols}
\noindent {\Large Zubereitung}\\
\textit{Vorbereitung:} Den Ofen auf \SI{200}{\celsius} Ober-/Unterhitze vorheizen.
Die Hähnchenbrust in Stücke schneiden.
Aus Paprikapulver, Limonen- bzw. Zitronensaft, Salz, Joghurt, Cayennepfeffer, Garam Masala, geriebenem bzw. kleingeschnittenem Ingwer und kleingeschnittenem bzw. gepresstem Knoblauch eine Marinade herstellen.
Das Fleisch mit der Marinade in einer Auflaufform vermischen und mindestens \SI{1}{h} einziehen lassen.
Noch besser ist es, das Fleisch bereits am Vortag zu marinieren und über Nacht in den Kühlschrank zu stellen.
Das Fleisch in der Auflaufform für \SI{25}{min} im vorgeheizten Ofen garen.
Die Zwiebel klein hacken und in \SI{2}{EL} Butter glasig anschwitzen.
Die passierten Tomaten, Zimt, Salz, Cayennepfeffer, geriebenen bzw. kleingeschnittenen Ingwer und kleingeschnittenen bzw. gepressten Knoblauch hinzugeben.
Alles \SI{20}{min} mit Deckel bei niedriger Temperatur köcheln lassen und gelegentlich umrühren.
Die restlichen \SI{2}{EL} Butter, den Honig und die Sahne hinzufügen und weitere \SI{3}{min} köcheln lassen.
Das Fleisch aus der Marinade nehmen und in die Soße geben.
Kurz umrühren und \SI{2}{min} mitköcheln lassen.
Dazu passen Reis oder Naan.
